% A self-contained layout example, not an official submission template.
% Change twocolumn to onecolumn to compare the same content at both widths.
\documentclass[twocolumn]{article}
\usepackage{amsmath,booktabs}
\pagestyle{empty}
\begin{document}
\section*{A layout example}
This example demonstrates local widths and cross-references without external
assets. The displayed values are illustrative, not measured results.

\section{Local widths}
The figure's panel occupies 80\% of the containing column. Inside the panel,
\verb|\linewidth| refers to the minipage's width. The same source therefore
adapts when the class option changes from two columns to one.

\begin{figure}[htbp]
  \centering
  \begin{minipage}{0.8\linewidth}
    \noindent\rule{\linewidth}{0.4pt}\par
    \smallskip
    \centering
    Column-sized panel\par
    The rules follow the local width.
    \par\smallskip
    \noindent\rule{\linewidth}{0.4pt}
  \end{minipage}
  \caption{An illustrative panel with no external graphics.}
  \label{fig:panel}
\end{figure}

\begin{table}[htbp]
  \centering
  \caption{Illustrative values with retained precision.}
  \label{tab:illustration}
  \begin{tabular*}{\linewidth}{@{\extracolsep{\fill}}lrr@{}}
    \toprule
    Method & Mean & Spread \\
    \midrule
    Alpha & 76.10 & 0.30 \\
    Beta & 78.20 & 0.40 \\
    \bottomrule
  \end{tabular*}
\end{table}

\section{References and math}
See Figure~\ref{fig:panel}, Table~\ref{tab:illustration}, and
equation~\eqref{eq:loss}.
\begin{equation}
  \mathcal{L}=\mathcal{L}_1+\lambda\mathcal{L}_2,
  \qquad \lambda=0.5.
  \label{eq:loss}
\end{equation}
The labels follow the counters they refer to. Build twice so the references
can resolve. Use the actual author kit for a venue submission.
\end{document}
